% ============================================================
%  MAPA MENTAL VIVO — Parcial I, Física 1
%
%  ARCHIVO DERIVADO. No editar a mano: lo regenera ENTERO el
%  agente generador-mapa-mental a partir de mapa-estudio.json
%  cada vez que se cierra una síntesis de un tema del Parcial I.
%  Cualquier edición manual se pierde en la próxima regeneración.
%
%  Estado al generar (2026-09-04):
%    - Temas del Parcial I: 12 · con síntesis aprobada: 3 · pendientes: 9
%      Cubiertos: Cinemática 1D (Guía 1), Dinámica I y Dinámica II (Guía 2).
%    - conceptos_ya_cubiertos (fisica-1): 33 — 28 en dificultad basica,
%      5 en intermedia. Las 5 intermedias tienen veces_reaparecido=1: son
%      las de Dinámica I que Dinámica II volvió a exigir. Se marcan con
%      $\star$, NO con color: el color codifica rama, y un color por
%      concepto individual rompería la lectura del agrupamiento.
%
%  ---------------------------------------------------------------
%  POR QUÉ AHORA SON DOS PANELES Y NO UN SOLO MAPA
%  ---------------------------------------------------------------
%  La vuelta anterior el mapa tenía 17 conceptos y entraba en un solo
%  tikzpicture. Con 33 ya NO entra, y no es un problema de agrupación:
%  es el techo de la plantilla. El cálculo de colisión de la vuelta
%  pasada, llevado a todos los pares de nodos (no solo a los extremos
%  de cada abanico), da esto con las constantes del preámbulo
%  (level distance 5.4 y 3.6; sibling angle 50; diámetros 3.6/2.7/2.1):
%
%    - Dos hojas hermanas de una misma rama quedan a 2*3.6*sin(25°) =
%      3.04 cm de centro a centro, o sea 0.94 cm de holgura entre bordes.
%      Ese 0.94 es el TECHO ABSOLUTO de holgura de cualquier reparto.
%    - Barriendo todas las particiones de N hojas en 4 a 11 ramas, con
%      2 a 6 hojas por rama, y todos los órdenes cíclicos de cada una:
%        N <= 21  ->  hay reparto con holgura positiva
%        N >= 22  ->  NINGÚN reparto evita el solapamiento
%      Para N = 33 el mejor caso posible es 4-2-3-4-4-4-4-4-4, y aun así
%      dos hojas se pisan 0.61 cm. No hay agrupación que lo salve.
%
%  El modelo se validó contra los dos casos documentados en la vuelta
%  anterior: reproduce +0.94 cm para el reparto 4-3-3-4-3 que compiló
%  bien, y -0.28 cm para el 4-3-2-4-4 que se pisaba, señalando el mismo
%  par de hojas entre las dos ramas de 4 contiguas.
%
%  Ante eso, y como las reglas prohíben tocar sibling angle, text width
%  o las distancias desde el cuerpo, la salida se parte en dos paneles
%  de una página cada uno, cortados por donde el contenido ya venía
%  cortado (Guía 1 / Guía 2). Cada panel conserva la relación 1 a 1
%  entre concepto del JSON y hoja del mapa: no se fusionó ni se omitió
%  ningún concepto. Ver el reporte al usuario: si se quiere volver a un
%  panel único habría que subir `level distance` del nivel 1 en el
%  preámbulo del plugin, que es un cambio aguas arriba, no local.
%
%  ---------------------------------------------------------------
%  PANEL 1 — Cinemática 1D + herramientas de cálculo (16 hojas, 6 ramas)
%  ---------------------------------------------------------------
%  Familias del JSON usadas como base del reparto, como corresponde
%  (esta es la primera vuelta en que el JSON trae familias finas):
%    "Cinematica 1D" (11) se subdivide en 4 ramas de lectura, porque 11
%    hojas en una sola rama es ilegible y supera el ~6 recomendado:
%      "Marco y magnitudes"       -> 3
%      "Derivadas del movimiento" -> 2
%      "MRU y MRUV"               -> 3
%      "Leer y resolver"          -> 3
%    "Calculo diferencial aplicado" (3) y "Calculo integral aplicado" (2)
%    entran tal cual, una rama cada una:
%      "Cálculo diferencial"      -> 3
%      "Cálculo integral"         -> 2
%
%  ORDEN CÍCLICO: 3-2-3-2-3-3, con las dos ramas de cálculo puestas una
%  al lado de la otra. Es la primera vez que el mapa incorpora contenido
%  que no es física sino herramienta matemática, y era la brecha crítica
%  declarada del lector, así que conviene que se distinga. El orden
%  además sigue el hilo del capítulo: describir -> derivar ->
%  [herramientas] -> integrar de vuelta a los modelos -> aplicar.
%  Holgura mínima verificada: +0.84 cm.
%
%  COLORES DEL PANEL 1: no se asignan en el orden de la paleta, sino
%  eligiendo qué color va en qué posición del ciclo. El primer intento
%  dejaba "Derivadas del movimiento" en ramaDos (marrón-gris) pegado a
%  "Cálculo diferencial" en ramaCinco (ocre): en el PDF los dos se
%  confundían, justo entre la rama de física y la de herramienta que más
%  importaba separar. Reasignado para que ningún vecino de una rama de
%  cálculo comparta tono con ella, sin mover ninguna hoja:
%      pos 1  Marco y magnitudes        ramaUno    azul-gris
%      pos 2  Derivadas del movimiento  ramaTres   verde-gris
%      pos 3  Cálculo diferencial       ramaCinco  ocre        <- cálculo
%      pos 4  Cálculo integral          ramaSeis   teal-gris   <- cálculo
%      pos 5  MRU y MRUV                ramaCuatro violeta-gris
%      pos 6  Leer y resolver           ramaDos    marrón-gris
%  El marrón queda en la posición 6, con vecinos violeta y azul, lejos
%  del ocre de la posición 3. Sigue habiendo un color por rama y ninguno
%  por concepto individual.
%
%  ---------------------------------------------------------------
%  PANEL 2 — Dinámica newtoniana (17 hojas, 5 ramas)
%  ---------------------------------------------------------------
%  Es exactamente el mapa de la vuelta anterior, sin tocar: mismo
%  reparto 4-3-3-4-3, mismos nombres y mismos colores de rama. Se
%  conserva a propósito, porque el lector ya viene asociando
%  azul=leyes y marrón=fuerzas, y romper esa asociación no aportaría
%  nada. La familia "Dinamica newtoniana" del JSON sigue trayendo los
%  17 juntos, así que la subdivisión en 5 ramas se mantiene igual.
%  Holgura mínima verificada: +0.94 cm (el techo de la plantilla).
%
%  NOTA sobre los colores: cada panel usa la paleta desde ramaUno, así
%  que un mismo color significa cosas distintas en un panel y en otro.
%  No genera ambigüedad porque son páginas separadas, cada una con su
%  nodo raíz y sus ramas rotuladas.
%
%  El mapa muestra SOLO conceptos con síntesis ya escrita. Los 9 temas
%  restantes del parcial no aparecen como nodos; su ausencia se declara
%  en el encabezado, que es texto de estado, no contenido inventado.
%
%  Ya NO aplica la restricción de la vuelta anterior de evitar notación
%  de derivada e integral: esta síntesis es justamente la que cerró esa
%  brecha, así que el panel 1 usa la notación con normalidad.
% ============================================================
\documentclass[11pt]{article}
% ============================================================
%  PLANTILLA DE MAPA MENTAL — syntheca
%  Clase STANDALONE/ARTICLE, usa la librería `mindmap` de TikZ
%  (distribución radial automática, no hay que calcular ángulos
%  a mano por nodo).
%
%  Filosofía de color, distinta al resto del sistema: acá SÍ se
%  usa color para diferenciar ramas (familias de conceptos) —
%  no es decoración, es una ayuda de codificación dual (dual
%  coding) para la memoria, la misma lógica por la que
%  pedagogia-cognitiva permite diagramas cuando cumplen una
%  función real. La paleta es apagada/profesional, nunca
%  saturada — mantiene el mismo espíritu sobrio del resto.
%
%  REGENERACIÓN COMPLETA: a diferencia del libro (que acumula
%  capítulos), este archivo se reescribe ENTERO cada vez que
%  generador-mapa-mental corre — nunca se edita a mano un nodo
%  suelto. La fuente de verdad es mapa-estudio.json, no el .tex.
%
%  ------------------------------------------------------------
%  COPIA LOCAL DEL REPO. El original vive en el plugin, en
%  ~/.claude/skills/syntheca/skills/plantilla-sintesis/latex/.
%  Esta copia es la que se edita; el plugin NO se toca.
%
%  CAMBIOS DE ESTA COPIA respecto de la del plugin (2026-09-05,
%  al pasar el mapa del Parcial I de 33 a 54 conceptos):
%    1. Tercer nivel de jerarquía (`level 3 concept`), para que
%       el mapa sea raíz -> familia -> subtema -> concepto y no
%       raíz -> familia -> concepto.
%    2. Distancias y anchos de texto mayores en todos los
%       niveles (era el fix ya identificado y nunca aplicado).
%    3. Lienzo configurable con \setLienzo: el diámetro natural
%       de un mapa radial crece de forma lineal con la cantidad
%       de hojas, así que el papel no puede ser fijo.
%    4. Se eliminó `\tikzstyle{every node}=[font=\sffamily]`.
%       Era un BUG: `every node` se aplica DESPUÉS de los
%       estilos de nivel, así que pisaba el `font=` de
%       `level 1/2/3 concept` y todos los nodos salían en
%       \normalsize. La familia sans ya la garantiza
%       \familydefault, así que borrarlo no cambia la tipografía
%       pero devuelve el control del CUERPO de letra por nivel,
%       que es de donde sale casi toda la capacidad de empaque.
%  ============================================================
\usepackage[utf8]{inputenc}
\usepackage{lmodern}   % antes de fontenc, para no caer en fuentes bitmap (Type3)
\usepackage[T1]{fontenc}
\usepackage{tikz}
\usetikzlibrary{mindmap,trees,backgrounds}
\usepackage{adjustbox}   % para encajar el mapa en la página (ver "Escalado")
\usepackage{fancyhdr}
\usepackage[table]{xcolor}
\usepackage{geometry}
\geometry{top=1.2cm, bottom=1.2cm, left=1cm, right=1cm, landscape}

% ---------- Lienzo ----------
% Un mapa radial tiene radio proporcional a la cantidad de hojas (todas las
% hojas viven sobre una misma circunferencia, así que el perímetro tiene que
% alcanzar para todas). Su ÁREA, entonces, crece con el cuadrado de la
% cantidad de conceptos: pasar de 17 hojas a 54 no triplica el mapa, lo
% multiplica por diez. Por eso el papel no puede ser fijo: con A4 apaisado
% un mapa de 54 conceptos sólo entra encogido a ~0.27, y a esa escala la
% letra queda en 2 pt. \setLienzo permite que cada mapa declare el papel que
% su cantidad de conceptos necesita, manteniendo el cuerpo de letra legible.
% Se llama en el cuerpo, antes de \begin{document}? No: geometry admite
% reconfiguración, así que se llama justo después de cargar el preámbulo.
\newcommand{\setLienzo}[2]{\geometry{paperwidth=#1, paperheight=#2, landscape=false,
  top=1.2cm, bottom=1.2cm, left=1.2cm, right=1.2cm}}

% Misma familia sans que el libro.
\renewcommand{\familydefault}{\sfdefault}

\definecolor{textgray}{RGB}{60, 60, 60}
\definecolor{rulegray}{RGB}{150, 150, 150}

% ---------- Paleta de ramas — apagada, profesional, máximo 6 colores.
% Si un examen tiene más de 6 familias, se reciclan en orden. ----------
\definecolor{ramaUno}{RGB}{91, 110, 130}    % azul-gris apagado
\definecolor{ramaDos}{RGB}{120, 100, 90}    % marrón-gris apagado
\definecolor{ramaTres}{RGB}{95, 120, 100}   % verde-gris apagado
\definecolor{ramaCuatro}{RGB}{115, 95, 120} % violeta-gris apagado
\definecolor{ramaCinco}{RGB}{130, 110, 80}  % ocre apagado
\definecolor{ramaSeis}{RGB}{90, 115, 120}   % teal-gris apagado

\newcommand{\miNombre}{Nombre Apellido}
\newcommand{\materiaActual}{Materia}
\newcommand{\setMateria}[1]{\renewcommand{\materiaActual}{#1}}
\newcommand{\nombreMapa}{Mapa mental}
\newcommand{\setNombreMapa}[1]{\renewcommand{\nombreMapa}{#1}}

\pagestyle{fancy}
\fancyhf{}
\renewcommand{\headrulewidth}{0pt}
\renewcommand{\footrulewidth}{0pt}
\fancyfoot[C]{\small\color{textgray}\thepage}
\fancyfoot[R]{\footnotesize\color{textgray}\miNombre}
\fancyfoot[L]{\footnotesize\color{textgray}\materiaActual}

% ---------- Estilo base de los conceptos (nodos) ----------
% Nivel raíz: el examen/tema central.
% Nivel 1: familias de conceptos (ramas, coloreadas).
% Nivel 2: subtemas dentro de una familia (heredan el color de la rama).
% Nivel 3: conceptos individuales (hojas, color heredado de su rama).
%
% El nivel 2 es NUEVO. Antes las hojas colgaban directo de la familia y una
% familia de 17 conceptos abría un abanico de 800°, imposible de acomodar.
% Con un nivel intermedio la misma familia se reparte en 2 o 3 abanicos
% chicos, y el ancho angular que necesita cae a menos de la mitad.
%
% grow cyclic es obligatorio en los tres niveles: sin él las ramas no se
% distribuyen en círculo, crecen todas hacia abajo y se apilan.
%
% text width fija el ancho de párrafo dentro del círculo; sin él, una
% etiqueta larga se desborda del nodo o lo infla de forma desproporcionada.
% minimum size va un poco por encima de text width para que quede aire.
%
% Las distancias de abajo NO son decorativas: salieron de una búsqueda sobre
% la grilla (level distance x sibling angle) verificando, para cada
% combinación, que no hubiera ni un par de nodos con distancia entre centros
% menor que la suma de sus radios. Cambiar una a ojo rompe esa garantía.
\tikzset{
  root concept/.style={
    concept color=textgray,
    font=\sffamily\bfseries\Large,
    text=white,
    minimum size=4.4cm,
    text width=3.6cm,
  },
  level 1 concept/.append style={
    grow cyclic,
    font=\sffamily\bfseries\small,
    text=white,
    level distance=19cm,
    minimum size=3.0cm,
    text width=2.5cm,
  },
  level 2 concept/.append style={
    grow cyclic,
    font=\sffamily\bfseries\footnotesize,
    text=white,
    sibling angle=52,
    level distance=7.5cm,
    minimum size=2.45cm,
    text width=2.0cm,
  },
  level 3 concept/.append style={
    grow cyclic,
    font=\sffamily\scriptsize,
    text=white,
    % La cuerda entre dos hojas hermanas es 2*d*sin(a/2); con d=4.6cm y a=34
    % da 2.69cm, contra un diámetro de hoja de ~2.13cm: sobran 0.56cm. Ese
    % número es la holgura mínima de TODO el mapa (ningún otro par queda más
    % cerca), así que bajar `sibling angle` o `level distance` acá es lo
    % primero que lo rompe.
    sibling angle=34,
    level distance=4.6cm,
    minimum size=2.0cm,
    text width=1.7cm,
  },
}

% ---------- Número de sectores angulares ----------
% El ángulo entre hijos de la raíz es 360/n. Con \setRamas{<familias>} todas
% las ramas reciben la MISMA porción, lo que sólo sirve si todas tienen una
% cantidad parecida de hojas. Cuando no la tienen (Parcial I: familias de 2 y
% de 17 conceptos), conviene repartir el círculo en SLOTS iguales y darle a
% cada rama varios slots, proporcionalmente a su cantidad de hojas: se llama
% \setRamas{<slots totales>} y en el cuerpo se rellena con `child[missing]`,
% que consume un slot sin dibujar nada. \setSlots es el mismo comando con el
% nombre que corresponde a ese uso.
\newcommand{\setRamas}[1]{%
  \tikzset{level 1 concept/.append style={sibling angle=360/#1}}%
}
\let\setSlots\setRamas
\setRamas{4}

% ---------- Escalado ----------
% Para que el mapa entre en la página, envolver TODO el tikzpicture en el
% entorno \begin{mapaAjustado} ... \end{mapaAjustado}, que lo encoge sólo si
% hace falta para respetar ancho Y alto disponibles. NUNCA usar scale=...
% junto con transform shape: grow cyclic rota el sistema de coordenadas de
% cada rama y transform shape le aplica esa rotación también al texto, así
% que la mitad de las etiquetas salen cabeza abajo.
%
% Ojo: `mapaAjustado` es una red de seguridad, no un plan. Si el mapa entra
% sólo gracias a un encogido fuerte, el problema es el lienzo (\setLienzo),
% porque el encogido se lleva puesto el cuerpo de letra.
\newenvironment{mapaAjustado}{%
  \begin{center}%
  \begin{adjustbox}{max size={\linewidth}{0.88\textheight}}%
}{%
  \end{adjustbox}%
  \end{center}%
}


\renewcommand{\miNombre}{Pedro Villar}
\setMateria{Física 1}
\setNombreMapa{Parcial I}

\begin{document}

% ============================================================
%  PANEL 1 — Cinemática 1D y herramientas de cálculo
% ============================================================
\thispagestyle{fancy}

\begin{center}
    {\Large\bfseries\color{textgray} Parcial I \textemdash{} Física 1\par}
    \vspace{0.10cm}
    {\color{rulegray}\rule{0.5\textwidth}{0.4pt}\par}
    \vspace{0.12cm}
    {\small\bfseries\color{textgray}
        Panel 1 de 2 \textemdash{} Cinemática 1D y herramientas de cálculo
        (Guía 1)\par}
    \vspace{0.05cm}
    {\footnotesize\color{rulegray}
        Mapa vivo al 4 de septiembre de 2026; el parcial es el 24 (faltan
        20 días). Sólo figura lo que ya tiene síntesis escrita: 33 conceptos
        de 3 de los 12 temas del parcial, y los otros 9 todavía no están
        cubiertos. Va en dos paneles porque 33 conceptos no entran en una
        sola estrella sin que las hojas se pisen. Las dos ramas en ocre y
        gris verdoso no son física: son el cálculo diferencial e integral
        que este capítulo trajo como herramienta.\par}
\end{center}

\vspace{0.10cm}

% La distribución radial la calcula sola la librería mindmap a partir de
% \setRamas{6} y de la lista de hijos. Acá NO se tocan grow cyclic,
% sibling angle ni text width: están fijados en el preámbulo. El grupo
% mantiene el \setRamas local a este panel, para que el panel 2 pueda
% pedir otra cantidad de ramas sin arrastrar ésta.
\begingroup
\setRamas{6}
\begin{mapaAjustado}
\begin{tikzpicture}[mindmap]
  \node[root concept, concept] {Cinemática 1D}
    child[concept color=ramaUno] { node[concept] {Marco y magnitudes}
      child { node[concept] {Sistema de referencia y $x(t)$} }
      child { node[concept] {Despla\-za\-mien\-to vs. camino} }
      child { node[concept] {Velocidad media y promedio} }
    }
    child[concept color=ramaTres] { node[concept] {Derivadas del movimiento}
      child { node[concept] {Velocidad: derivada de $x(t)$} }
      child { node[concept] {Ace\-le\-ra\-ción: segunda derivada} }
    }
    child[concept color=ramaCinco] { node[concept] {Cálculo diferencial}
      child { node[concept] {Límite y cociente incremental} }
      child { node[concept] {Reglas de derivación} }
      child { node[concept] {Críticos, concavidad, inflexión} }
    }
    child[concept color=ramaSeis] { node[concept] {Cálculo integral}
      child { node[concept] {Primitiva y constante $C$} }
      child { node[concept] {Barrow: la integral es área} }
    }
    child[concept color=ramaCuatro] { node[concept] {MRU y MRUV}
      child { node[concept] {Deducidos por integración} }
      child { node[concept] {Relación sin tiempo} }
      child { node[concept] {Caída libre y tiro vertical} }
    }
    child[concept color=ramaDos] { node[concept] {Leer y resolver}
      child { node[concept] {Tramos y continuidad de $v$} }
      child { node[concept] {Encuentro de dos móviles} }
      child { node[concept] {Gráficos $x$, $v$, $a$ juntos} }
    };
\end{tikzpicture}
\end{mapaAjustado}
\endgroup

\newpage

% ============================================================
%  PANEL 2 — Dinámica newtoniana
% ============================================================
\thispagestyle{fancy}

\begin{center}
    {\Large\bfseries\color{textgray} Parcial I \textemdash{} Física 1\par}
    \vspace{0.10cm}
    {\color{rulegray}\rule{0.5\textwidth}{0.4pt}\par}
    \vspace{0.12cm}
    {\small\bfseries\color{textgray}
        Panel 2 de 2 \textemdash{} Dinámica newtoniana (Guía 2)\par}
    \vspace{0.05cm}
    {\footnotesize\color{rulegray}
        $\star$ marca los 5 conceptos que Dinámica II volvió a exigir a
        mayor dificultad: son el núcleo ya consolidado.\par}
\end{center}

\vspace{0.10cm}

\begingroup
\setRamas{5}
\begin{mapaAjustado}
\begin{tikzpicture}[mindmap]
  \node[root concept, concept] {Dinámica}
    child[concept color=ramaUno] { node[concept] {Las tres leyes}
      child { node[concept] {Fuerza: causa del cambio} }
      child { node[concept] {Primera ley: marco inercial} }
      child { node[concept] {$\star$ Segunda ley, eje por eje} }
      child { node[concept] {$\star$ Tercera ley: pares} }
    }
    child[concept color=ramaDos] { node[concept] {Masa y tipos de fuerza}
      child { node[concept] {Masa inercial} }
      child { node[concept] {$\star$ Peso, normal y tensión} }
      child { node[concept] {Ley de Hooke} }
    }
    child[concept color=ramaTres] { node[concept] {Cómo se plantea}
      child { node[concept] {$\star$ Diagrama de cuerpo aislado} }
      child { node[concept] {$\star$ Plano inclinado} }
      child { node[concept] {Signos, cuerpo por cuerpo} }
    }
    child[concept color=ramaCuatro] { node[concept] {Roza\-mien\-to}
      child { node[concept] {Estático: una desigualdad} }
      child { node[concept] {Dinámico y coeficientes} }
      child { node[concept] {Opuesto al deslizamiento} }
      child { node[concept] {Ángulo crítico mide $\mu_e$} }
    }
    child[concept color=ramaCinco] { node[concept] {Sistemas acoplados}
      child { node[concept] {Vínculo cinemático} }
      child { node[concept] {Atajo del conjunto y polea ideal} }
      child { node[concept] {Roza\-mien\-to y normales apiladas} }
    };
\end{tikzpicture}
\end{mapaAjustado}
\endgroup

\end{document}
